\documentclass[11pt]{article}
\usepackage[T1]{fontenc}
\usepackage{textcomp}
\usepackage[margin=0.75in]{geometry}
\usepackage{amsmath,amssymb,amsthm}
\usepackage{array,booktabs}
\usepackage{hyperref}
\setlength{\parindent}{0pt}
\newtheorem{theorem}{Theorem}
\theoremstyle{definition}
\newtheorem{definition}{Definition}
\title{Flow Proof Artifact: prop-06-equal-angles-equal-sides}
\author{Generated by \texttt{flow doc proof}}
\date{Flow Proof Book}
\begin{document}
\maketitle
If two angles of a triangle are equal, the sides opposite them are equal.\\[0.5em]
\textbf{Source.} euclid — Elements, Book I, Proposition 6\\[0.5em]
\bigskip
\noindent{\large \textbf{Derived fact 1.} \textit{If two angles of a triangle are equal, the sides opposite them are equal} \hfill the Euclidean plane \cdot Euclid Book I \cdot Proposition 6: equal angles imply equal opposite sides}\par
\smallskip\noindent\textit{Coordinate: the Euclidean plane \cdot Euclid Book I \cdot Proposition 6: equal angles imply equal opposite sides}\par
\smallskip\noindent\textit{Source: euclid — Elements, Book I, Proposition 6}\par
\smallskip\noindent\textit{Built on: proposition 5: base angles of an isosceles triangle are equal, for Book I of the Elements in on the Euclidean plane}\par
\medskip
\noindent\textbf{Goal.} If two angles of a triangle are equal, the sides opposite them are equal\par
\noindent$\displaystyle \text{side AC} = \text{side AB}$\par
\medskip
\noindent
\begin{tabular}{@{} >{\raggedright\arraybackslash}p{0.06\textwidth} >{\raggedright\arraybackslash}p{0.47\textwidth} >{\raggedleft\arraybackslash}p{0.41\textwidth} @{}}
\textbf{\#} & \textbf{Proof} & \textbf{Mathematics} \\
\midrule
\textbf{1.} & We prove that proposition 6: equal angles imply equal opposite sides for Book I of the Elements in the Euclidean plane. & \\[0.45em]
\textbf{2.} & We invoke the derived fact: angle\_ABC == angle\_ACB. & \\[0.45em]
\textbf{3.} & We invoke the derived fact governing Book I of the Elements in the Euclidean plane: proposition 5: base angles of an isosceles triangle are equal, for Book I of the Elements in on the Euclidean plane. & \\[0.45em]
\textbf{4.} & From step 2 and step 3, we can deduce that segment AB equals segment AC. & $\displaystyle \text{segment AB} = \text{segment AC}$ \\[0.45em]
\textbf{5.} & We can deduce that side AC equals side AB. Hence proven. & $\displaystyle \text{side AC} = \text{side AB}$ \\[0.45em]
\bottomrule
\end{tabular}
\label{thm:Geometry-Euclid Book I-proposition_6_equal_angles_imply_equal_opposite_sides}
\par\medskip\hrule
\end{document}
